% Part: turing-machines
% Chapter: undecidability
% Section: representing-tms

\documentclass[../../../include/open-logic-section]{subfiles}

\begin{document}

\olfileid[fr]{tur}{und}{rep}
\olsection{Représenter les machines de Turing}

\begin{explain}
Pour représenter les machines de Turing et leur comportement par
!!a{sentence} de logique du premier ordre, il nous faut définir un
langage approprié. Celui-ci comporte deux sortes de ressources :
des !!{predicate}s décrivant les configurations de la machine et
des expressions permettant de numéroter les étapes du calcul
(les « instants ») et les positions sur le ruban.

Nous introduisons deux familles de !!{predicate}s binaires :
pour chaque état~$q$, !!a{predicate}~$\Obj Q_q$; pour chaque
symbole du ruban~$\sigma$, !!a{predicate}~$\Obj S_\sigma$.
Les premiers servent à décrire l'état de~$M$ et la position de sa
tête; les seconds, le contenu du ruban.

Pour exprimer les positions de la tête et le nombre d'étapes
exécutées, nous devons pouvoir représenter des nombres. Nous
utilisons à cette fin !!a{constant}~$\Obj 0$ et un symbole de fonction
unaire~$\prime$, la fonction successeur. Par convention, ce symbole
s'écrit \emph{après} son argument, sans parenthèses.

À chaque nombre~$n$ correspond un terme canonique~$\num{n}$,
le \emph{numéral} de~$n$, qui le représente dans~$\Lang L_M$.
Ainsi, $\num{0}$ est $\Obj 0$, $\num{1}$ est $\Obj 0'$,
$\num{2}$ est $\Obj 0''$, et ainsi de suite. Plus précisément :
\begin{align*}
\num{0} & = \Obj 0 \\
\num{n+1} &= \num{n}'
\end{align*}
Le terme~$\num{0}$, c'est-à-dire $\Obj 0$, désigne à la fois la
case la plus à gauche du ruban et l'instant qui précède la première
étape du calcul : celui de la configuration initiale. Le terme
$\num{1}$, c'est-à-dire $\Obj 0'$, désigne à la fois la case située
immédiatement à droite de cette première case et l'instant qui suit
la première étape; et ainsi de suite.

Nous introduisons également !!a{predicate}~$<$ pour exprimer
l'ordre des positions sur le ruban (« à gauche de ») et celui des
étapes du calcul (« avant »).

Le langage étant fixé, nous pouvons énumérer les « axiomes » de
$!T(M,w)$ : ce sont les !!{sentence}s qui, pris ensemble, décrivent
le comportement de~$M$ sur l'entrée~$w$. Certains !!{sentence}s
portent sur $\Obj 0$, $\prime$ et~$<$; d'autres !!{sentence}s
décrivent la configuration initiale; les derniers !!{sentence}s
décrivent la configuration qui résulte de l'exécution de chaque
instruction.
\end{explain}

\begin{defn}
\ollabel{defn:tm-descr}
Étant donné une machine de Turing $M=\tuple{Q,\Sigma,q_0,\delta}$,
le langage~$\Lang L_M$ comprend :
\begin{enumerate}
\item Pour chaque état $q\in Q$, un !!{predicate} binaire
  $\Obj Q_q(x,y)$. Intuitivement, $\Obj Q_q(\num{m},\num{n})$
  signifie qu'après $n$ étapes, $M$ est dans l'état~$q$ et sa tête
  lit la case d'indice~$m$ (la case la plus à gauche a l'indice~$0$).
\item Pour chaque symbole $\sigma\in\Sigma$, un !!{predicate}
  binaire $\Obj S_\sigma(x,y)$. Intuitivement,
  $\Obj S_\sigma(\num{m},\num{n})$ signifie qu'après $n$ étapes,
  la case d'indice~$m$ contient le symbole~$\sigma$.
\item Une !!{constant} $\Obj 0$.
\item Un !!{function} unaire $\prime$.
\item Un !!{predicate} binaire $<$.
\end{enumerate}
\end{defn}

Les !!{sentence}s qui décrivent le fonctionnement de la machine
de Turing~$M$ sur l'entrée $w=\sigma_{i_1}\dots\sigma_{i_k}$
sont les suivants :
\begin{enumerate}
\item Axiomes portant sur les nombres et sur~$<$ :
\begin{enumerate}
\item !!^a{sentence} affirmant que tout nombre est inférieur à son
  successeur :
\[
\lforall[x][x<x']
\]
\item !!^a{sentence} assurant la transitivité de~$<$ :
\[
\lforall[x][\lforall[y][\lforall[z][
  ((x<y\land y<z)\lif x<z)]]]
\]
\end{enumerate}
\item Axiomes décrivant la configuration initiale :
\begin{enumerate}
\item Avant la première étape, la machine~$M$ est dans l'état
  initial~$q_0$ et sa tête lit la case~$1$ :
\[
\Obj Q_{q_0}(\num{1},\num{0})
\]
\item Les $k+1$ premières cases contiennent respectivement les
  symboles $\TMendtape$, $\sigma_{i_1}$, \dots,
  $\sigma_{i_k}$ :
\[
\Obj S_\TMendtape(\num{0},\num{0})\land
\Obj S_{\sigma_{i_1}}(\num{1},\num{0})\land
\dots\land
\Obj S_{\sigma_{i_k}}(\num{k},\num{0})
\]
  Si $k=0$, cette conjonction se réduit au premier terme.
\item Toutes les autres cases sont blanches :
\[
\lforall[x][(\num{k}<x\lif\Obj S_\TMblank(x,\num{0}))]
\]
\end{enumerate}
\item Axiomes décrivant le passage d'une configuration à la
  suivante :

Dans ce qui suit, $!A(x,y)$ désigne la conjonction, pour tous les
symboles $\sigma\in\Sigma$, des !!{formula}s de la forme
\[
\lforall[z][
  (((z<x\lor x<z)\land\Obj S_\sigma(z,y))
  \lif\Obj S_\sigma(z,y'))]
\]
Les variables $x$ et~$y$ restent libres dans~$!A(x,y)$.\footnote{Correction éditoriale : la source appelle ces conjoints des énoncés ; ce sont des formules ouvertes, qui deviennent closes lorsque leurs variables libres sont liées dans les axiomes de transition.}
Nous employons $!A(\num{m},\num{n})$ pour exprimer qu'après $n+1$
étapes, le contenu du ruban est le même qu'après $n$ étapes,
sauf éventuellement dans la case~$m$.
\begin{enumerate}
\item \ollabel{rep-right} Pour chaque instruction
  $\delta(q_i,\sigma)=\tuple{q_j,\sigma',\TMright}$,
  l'!!{sentence} :
\begin{align*}
& \lforall[x][\lforall[y][(
  (\Obj Q_{q_i}(x,y)\land\Obj S_\sigma(x,y))\lif {}]]\\
&\qquad (\Obj Q_{q_j}(x',y')\land\Obj S_{\sigma'}(x,y')
  \land !A(x,y)))
\end{align*}
Si, après $y$ étapes, la machine est dans l'état~$q_i$ et sa tête
lit la case~$x$ contenant~$\sigma$, alors, une étape plus tard,
elle est dans l'état~$q_j$ et lit la case~$x+1$; la case~$x$
contient désormais~$\sigma'$, tandis que toutes les autres cases
conservent leur contenu.

\item \ollabel{rep-left} Pour chaque instruction
  $\delta(q_i,\sigma)=\tuple{q_j,\sigma',\TMleft}$,
  l'!!{sentence} :
\begin{align*}
& \lforall[x][\lforall[y][(
  (\Obj Q_{q_i}(x',y)\land\Obj S_\sigma(x',y))\lif {}]]\\
&\qquad (\Obj Q_{q_j}(x,y')\land\Obj S_{\sigma'}(x',y')
  \land !A(x',y)))\land {}\\
& \lforall[y][(
  (\Obj Q_{q_i}(\num{0},y)\land\Obj S_\sigma(\num{0},y))\lif {}]\\
&\qquad (\Obj Q_{q_j}(\num{0},y')
  \land\Obj S_{\sigma'}(\num{0},y')
  \land !A(\num{0},y)))
\end{align*}
Examinons ce fonctionnement. Cette fois, le point de départ n'est
pas « la tête lit la case~$x$ », mais « la tête lit la case~$x+1$ ».
Après un déplacement vers la gauche, elle lit la case~$x$, alors
que l'écriture a eu lieu dans la case~$x+1$. Cette formulation évite
d'avoir à introduire une soustraction ou une fonction prédécesseur.
Le terme de conservation est donc $!A(x',y)$ : il exclut de la
conservation la case~$x+1$, celle où l'instruction écrit.\footnote{Correction éditoriale : la source emploie ici $!A(x,y)$, ce qui conserverait l'ancien symbole de la case écrite $x+1$. Nous remplaçons ce terme par $!A(x',y)$ dans cette clause et dans sa vérification.}

Les nombres de la forme $x+1$ sont $1$, $2$, \dots{} : le premier
conjoint ne couvre donc pas le cas où la tête lit la case~$0$ et
reçoit l'instruction d'aller à gauche. Elle ne peut alors reculer
et reste en place. Le second conjoint traite ce cas : si, après
$y$ étapes, la tête lit la case~$0$ dans l'état~$q_i$ et que cette
case contient~$\sigma$, alors, une étape plus tard, la tête lit
encore la case~$0$, l'état est~$q_j$ et cette case contient
$\sigma'$; les autres cases conservent leur contenu.

\item \ollabel{rep-stay} Pour chaque instruction
  $\delta(q_i,\sigma)=\tuple{q_j,\sigma',\TMstay}$,
  l'!!{sentence} :
\begin{align*}
& \lforall[x][\lforall[y][(
  (\Obj Q_{q_i}(x,y)\land\Obj S_\sigma(x,y))\lif {}]]\\
&\qquad (\Obj Q_{q_j}(x,y')\land\Obj S_{\sigma'}(x,y')
  \land !A(x,y)))
\end{align*}
\end{enumerate}
\end{enumerate}
Notons~$!T(M,w)$ la conjonction de tous les !!{sentence}s qui
précèdent, pour la machine de Turing~$M$ et l'entrée~$w$.

Pour exprimer que~$M$ finit par s'arrêter, cherchons !!a{sentence}
affirmant qu'après un certain nombre d'étapes, la fonction de
transition n'est pas définie. Soit~$X$ l'ensemble des couples
$\tuple{q,\sigma}$ pour lesquels $\delta(q,\sigma)$ n'est pas
définie. Notons alors~$!E(M,w)$ l'!!{sentence}
\[
\lexists[x][\lexists[y][(\bigvee_{\tuple{q,\sigma}\in X}
  (\Obj Q_q(x,y)\land\Obj S_\sigma(x,y)))]]
\]
La disjonction vide, si $X=\varnothing$, vaut $\bot$.

Si nous employons une machine de Turing dotée d'un état d'arrêt
désigné~$h$, dans lequel tout arrêt a lieu et dont aucune instruction ne part, la définition se simplifie : l'!!{sentence}~$!E(M,w)$
\[
\lexists[x][\lexists[y][\Obj Q_h(x,y)]]
\]
exprime que la machine finit par s'arrêter.

\begin{prop}
\ollabel{prop:mlessk}
Si $m<k$, alors $!T(M,w)\Entails\num{m}<\num{k}$.
\end{prop}

\begin{proof}
La preuve est laissée en exercice.
\end{proof}

\begin{prob}
Démontrez la \olref[tur][und][rep]{prop:mlessk}.
Indication : raisonnez par récurrence sur~$k-m$.
\end{prob}

\end{document}
